\documentclass[11pt, a4paper, oneside]{article}

\usepackage{indentfirst}

% hifenização e outras especificações para português
\usepackage[portuguese]{babel}
\usepackage{listings}
% para ter imagens, depois define a directoria de imagens
\usepackage{graphicx}
\graphicspath{{./imagens/}}
\usepackage[utf8]{inputenc}    

\usepackage{listings}
\usepackage{color}
 
\definecolor{dkgreen}{rgb}{0,0.6,0}
\definecolor{gray}{rgb}{0.5,0.5,0.5}
\definecolor{mauve}{rgb}{0.58,0,0.82}
\usepackage{color}
\usepackage{listings}
\lstset{ %
language=C++,                % choose the language of the code
basicstyle=\footnotesize,       % the size of the fonts that are used for the code
numbers=left,                   % where to put the line-numbers
numberstyle=\footnotesize,      % the size of the fonts that are used for the line-numbers
numberstyle=\tiny,
stepnumber=1,                   % the step between two line-numbers. If it is 1 each line will be numbered
numbersep=5pt,                  % how far the line-numbers are from the code
keywordstyle=\color[rgb]{0,0,1},
basicstyle=\tiny,
commentstyle=\color[rgb]{0.133,0.545,0.133},
stringstyle=\color[rgb]{0.627,0.126,0.941},
showspaces=false,               % show spaces adding particular underscores
showstringspaces=false,         % underline spaces within strings
showtabs=false,                 % show tabs within strings adding particular underscores
frame=single,           % adds a frame around the code
tabsize=2,          % sets default tabsize to 2 spaces
captionpos=b,           % sets the caption-position to bottom
breaklines=true,        % sets automatic line breaking
breakatwhitespace=false,    % sets if automatic breaks should only happen at whitespace
escapeinside={\%*}{*)}          % if you want to add a comment within your code
}

% hiperligações
\usepackage{hyperref}
\hypersetup{colorlinks=true, urlcolor=blue, linkcolor=black}

% escrever acentos e coisas do género sem que o latex se desoriente
\usepackage[utf8]{inputenc}

% para ter imagens, depois define a directoria de imagens
\usepackage{graphicx}
\graphicspath{{./imagens/}}

\usepackage[labelformat=simple]{caption}
\usepackage[labelformat=empty]{subcaption}

% para ter a informação de quantas páginas tem o documento
\usepackage{lastpage}

% definir o cabeçalho e rodapé
\usepackage{fancyhdr}
\pagestyle{fancy}
\fancyhead[L]{\small{\textbf{Gestão de Campeonatos de Orientação}}}
\fancyhead[R]{\small{Processamento de Linguagens}}

% ter enumerações alinhadas
\usepackage{enumitem}

% escrever algoritmos
\usepackage[algoruled]{algorithm2e}

% definir comandos especiais
\newcommand\doubleplus{+\kern-1.3ex+\kern0.8ex} %

\newcommand{\todo}[1] {\textcolor{BrickRed}{\begin{quote}#1\end{quote}}}

%\usepackage{listings}


%%%%%%%%%%%%%%%%%%%%%%%%%%%%%%%%%%%%%%%%%%%%%%
%% inicio do documento
\begin{document}
\title{Gestão de Campeonatos de Orientação}
\date{\today\\Universidade do Minho}
\author{
  Bruno Ferreira\\
  {\small A61055}\\
  \and
  Cláudia Oliveira\\
  {\small A60987}\\
  \and
  Vanessa Campos\\
  {\small A54801}\\
}

\maketitle

\begin{figure}[h]
\begin{center}
\includegraphics[width=0.4\linewidth]{logo}
\end{center}
\end{figure}


\begin{abstract}

  O presente trabalho foi desenvolvido no âmbito da Unidade Curricular de Processamento de Linguagens e tem como principal objetivo o aumento da experiência na Linguagem Imperativa C, a capacidade de escrever Gramáticas Independentes de Contexto que satisfaçam a condição de LR(), a utilização de compiladores como o Lex/Yacc e o desenvolvimento de Processadores de Linguagem. De entre os variados temas, Gestão de Campeonatos de Orientação foi o tema escolhido pelo nosso grupo, para o desenvolvimento deste segundo trabalho prático. Ao longo deste relatório iremos explicar as Estruturas de Dados que foram implementadas para o desenvolvimento deste trabalho, bem como todas as decisões que foram tomadas ao longo do desenvolvimento deste projeto. 

\end{abstract}
\newpage

\tableofcontents
\listoffigures 

\newpage
\section{Introdução}

A Orientação é um desporto que tem como objetivo percorrer uma determinada distância, onde o atleta tem de obrigatoriamente passar pelos pontos que estão indicados no mapa que lhe é atribuído.

Neste contexto, o foco deste projeto passa pela gestão de campeonatos de orientação onde é necessário saber a informação das provas que são efetuadas, os participantes envolvidos e os tempos que cada atleta obteve, de modo a serem obtidas as classificações necessárias.

Um objetivo principal do projeto é a leitura de ficheiros onde estes tem uma configuração definida pelo ficheiro de configuração, depois a leitura do CSV com a informação dos jogadores e mais tarde o geramento de um página HTML onde é apresentada a informação tratada.

Para a implementação foi necessário a construção de gramáticas independentes de contexto, implementação de parses para cada ficheiro que necessário, e a utilização do lex e yacc para retirar a informação necessária.


\newpage
\section{Desenvolvimento}

\subsection{Contextualização do Problema}

Flex é uma ferramenta para gerar automaticamente analisadores léxicos, isto é, programas que reconhecem padrões léxicos num texto. O Flex é uma evolução da ferramenta Lex, mas com a caraterística de ser mais rápido que este (Fast Lex).

Yacc é utilizado para sistemas operacionais Unix, onde o seu principal objetivo e gerar analisadores sintáticos. No entanto o yacc tem de ser utilizado em conjunto com o lex, pois o yacc não consegue ler a partir de uma entrada de dados, daí a utilização  do lex para que este gere os tokens que mais tarde são usados pelo yacc para o processamento.

Gramática Independente de Contexto é uma gramática formal onde são definidas regras de produção da 
 V $\rightarrow$ W onde V é um símbolo não terminal e w um conjunto de símbolos terminais ou variáveis (W no entanto também pode ser nulo).

\subsection{Gestão de Campeonatos de Orientação}

Neste segundo trabalho de Processamento de Linguagens foi proposto a seleção de um projeto entre vários propostos. Após uma análise e chegada a uma concordância, o nosso grupo optou por escolher o tema de Gestão de Campeonato de Orientação.

Com este enunciado existem várias operações que ocorrem em simultâneo, e por isso existem vários ficheiros, como:
\begin{itemize}
\item Ficheiro de Resultados é um ficheiro CSV, onde são guardados os resultados obtidos em cada prova.

\item Ficheiro de Configuração é um ficheiro que vai definir o processamento a realizar, isto é quais as provas e os campos que vão ser tratados.
\end{itemize}

\subsection{Enunciado}

De uma maneira geral, para este trabalho pretende-se desenvolver os seguintes pontos:

\begin{itemize}
\item Especificação de uma gramática para a linguagem de comandos da consola;
\item Especificação de uma gramática para a linguagem de especificação de connfigurações;
\item Especificação de uma gramática para os ficheiros de resultados;
\item Desenvolver os respetivos parses;
\item Adicionar as ações semânticas necessárias,
\end{itemize}

\subsection{Descrição do Problema}

Utilizando o flex, yac e gramáticas independentes de contexto era necessário a desenvolvimento de uma aplicação, que permita o tratar ficheiros específicos e gere o HTML como os resultados pretendidos.

No entanto existem outros requisitos que permitem o bom funcionamento da aplicação. Deste modo, a aplicação deve possuir um ambiente de consola e implementar os seguintes comandos:
\begin{itemize}
\item \textbf{Carregar configuração} - Permite carregar o ficheiro de configuração;
\item \textbf{Carregar base de dados} - Carregar um ficheiro de dados previamente guardado;
\item \textbf{Carregar resultado de prova} - Carregar o ficheiro de resultados;
\item \textbf{Calcular Ranking} - Permite calcular o ranking de pontuação dos atletas;
\item \textbf{Gravar base de dados} - Guardar a informação que existem na memória num ficheiro;
\item \textbf{Sair} - Termina a aplicação.
\end{itemize}

É a partir daqui que os ficheiro de resultados e configuração são lidos e depois as operações são executadas. 

Antes de ser criado o HTML a aplicação tem de ser configurada.
A informação para tal é lida de um ficheiro de configuração onde nele são definidos os campos:
\begin{itemize}
\item \textbf{Titulo} - Titulo da prova em causa;
\item \textbf{nporvas} - Nº de provas a realizar;
\item \textbf{N} - Indicados do topN ;
\item \textbf{campos} - Ordem dos campos quando no CSV;
\item \textbf{score} - Pontuação obtida.
\end{itemize}

\newpage

\section{Desenvolvimento do Programa}
Um dos pressupostos da implementação deste projeto é o tratamento de diferente tipo de informação, foi decidido estruturar-se o projeto consoante as funções que cada ficheiro implica.
[INSERIR IMAGEM DOS FICHEIROS]

O desenvolvimento deste projeto passa assim a desenvolver três gramáticas independentes de contexto distintas: 
\begin{itemize}
\item Gramática para a linguagem de comandos da consola; 
\item Gramática para a linguagem de configurações;
\item Grámatica para ficheiros de resultados. 
\end{itemize}

Para a criação das estruturas de dados foi usada a ferramenta \textit{gabs} disponibilizada pelo professor que cria a estrutura e os seus construtores quando é executada com uma gramática abstrata como argumento.

De modo a utiliza-se esta ferramenta era necessário, primeiramente, a criação de uma gramática abstrata. Por isso, uma das primeiras etapas deste projeto foi desenvolver a especificação de cada uma gramáticas independentes do contexto. 

Em cada uma das subsecções seguintes, será apresentada cada uma das gramáticas desenvolvidas, bem como a sua estrutura de dados implementada, as assinaturas das funções em C e os seus construtores.

\subsection{Especificação de Configurações da Aplicação}
O ficheiro de configurações define o processamento a realizar: campos a extrair para o resultado final, o número de provas que pontuam, o número de melhores resultados que serão contabilizados para a classificação, os campos que irão aparecer na página de resultado.

%%%%%%%%%%%%%%%%%%%%%%%%%%GRAMÁTICA-CONFIGUraÇÕES%%%%%%%%%%%%%%%%%%%%%%%%%%%%%

\subsubsection{Gramática para a Linguagem de Especificação de Configurações}

A gramática a apresentar é a usada para o ficheiro de configuração. Assim sendo tem-se:

\begin{lstlisting}[language=C, caption={Gramática para a linguagem de especificação de
 configurações.}]
<ga>
  Confs -> cons_cfg_Confs(Conf Confs)
        |  cons_cfg_Confs_NIL()
        ;

  Conf -> cons_cfg_Conf_Titulo(STR)
       |  cons_cfg_Conf_Nprovas(INT)
       |  cons_cfg_Conf_Ntop(INT)
       |  cons_cfg_Conf_Campos(Lcampos)
       |  cons_cfg_Conf_Tempo(INT)
       |  cons_cfg_Conf_Chave(INT)
       |  cons_cfg_Conf_Nome(INT)
       ;

  Lcampos -> cons_cfg_Lcampos_Lcampos(Lcampos INT)
          |  cons_cfg_Lcampos_Campo(INT)
          ;
</ga>
\end{lstlisting}

\newpage
%%%%%%%%%%%%%%%%%%%%%%%%%%ESTRUTURAS-CONFIGUraÇÕES%%%%%%%%%%%%%%%%%%%%%%%%%%%%%
\subsubsection{Estrutura de Dados da Especificação de Configurações}
A estrutura que é gerada para o armazenamento dos campos que a gramática abstrata definiu, é divida consoante o nível em que se esta. Assim, começa-se pela estrutura que vai abragir mais campos, neste caso o total de configurações, depois pela configuração, e, por último, a estrutura que vai guardar um campo para uma configuração.

\begin{lstlisting}[language=C, caption={Estrutura de dados para armazenar as configurações.}]
struct sConfs 
{ int flag;
  union {  
    struct {
        Conf s1;
        Confs s2;
      } d1;
    struct {
      } d2;

  } u;
};
\end{lstlisting}

\begin{lstlisting}[language=C, caption={Estrutura de dados que armazena uma configuração.}]
struct sConf 
{ int flag;
  union {  
    struct {
        char * s1;
      } d1;
    struct {
        int s1;
      } d2;
    struct {
        int s1;
      } d3;
    struct {
        Lcampos s1;
      } d4;
    struct {
        int s1;
      } d5;

  } u;
};
\end{lstlisting}

\begin{lstlisting}[language=C, caption={Estrututura de dados que armazena os campos de um configuração.}]
struct sLcampos 
{ int flag;
  union {  
    struct {
        Lcampos s1;
        int s2;
      } d1;
    struct {
        int s1;
      } d2;

  } u;
};
\end{lstlisting}

Verifica-se aqui que cada \textit{struct} acima apresentada deriva de cada uma das produções que foram indicadas em na gramática abstrata.
As produções que foram usadas para definir as produções da gramática abstrata são as que também foram usadas para definir as produções do yacc.
\newpage

\subsubsection{Assinatura das funções utilizadas}
Aqui serão apresentadas as assinaturas das funções que foram utilizadas para nós dar suporte à implementação do nosso projecto.

\begin{lstlisting}[language=C, caption={Construtores do ficheiro de configuração.}]
/* -----------------------------------
 * Constructor Function Signatures
 * -----------------------------------
 */

Confs  cons_cfg_Confs( Conf a1, Confs a2);
Confs  cons_cfg_Confs_NIL();

Conf  cons_cfg_Conf_Titulo( char * a1);
Conf  cons_cfg_Conf_Nprovas( int a1);
Conf  cons_cfg_Conf_Ntop( int a1);
Conf  cons_cfg_Conf_Campos( Lcampos a1);
Conf  cons_cfg_Conf_Tempo( int a1);
Conf  cons_cfg_Conf_Chave( int a1);
Conf  cons_cfg_Conf_Nome( int a1);

Lcampos  cons_cfg_Lcampos_Lcampos( Lcampos a1, int a2);
Lcampos  cons_cfg_Lcampos_Campo( int a1);
\end{lstlisting} 

Os construtores da aplicação são criados automaticamente, isto é, é passado ao \textit{gabs} a gramática abstrata criada. Esta ferramenta para além de criar as estrutura de dados, cria também todos os construtores necessários para a criação da mesma.

\begin{lstlisting}[language=C, caption={Funções do ficheiro de configuração.}]
/* -----------------------------------
 * Custom Function Signatures
 * -----------------------------------
 */
void cfg_Confs_print( Confs cfgs );
int cfg_Confs_validate( Confs cfgs );
Lcampos cfg_Lcampos_reverse( Lcampos l );

// getters
int cfg_get_Nprovas( Confs cfgs );
int cfg_get_Ntop( Confs cfgs );
int cfg_get_Tempo( Confs cfgs );
int cfg_get_Chave( Confs cfgs );
int cfg_get_Nome( Confs cfgs );
char* cfg_get_Titulo( Confs cfgs );

// obtem um array de inteiros com os campos. terminado com -1
int* cfg_get_Campos( Confs cfgs );
int cfg_get_NCampos( Confs cfgs );

// obtem um array de bytes que esta a 1 nos campos que foram selecionados
char* cfg_Campos_seleccionado( Confs cfgs, int totalCampos );
\end{lstlisting} 

Embora já estejam os construtores definidos automaticamente era necessário criar novas funções para ser possível manusear a estrutura de dados. Deste modo, foram criadas funções que permitem ir "buscar" algum campo específico, ou fazer operações como de validação, imprimir, entre outras.

Como vai ser necessário aceder à informação que existe, definimos as funções que nos permitem dar os valores que são os campos de configurações.

\textbf{Nota:} Foi necessário desenvolver a função \texttt{Lcampos cfg\_Lcampos\_reverse( Lcampos l )}, uma vez que a gramática definida neste programa tem recursividade à esquerda e era necessário imprimir os valores pela ordem correta, daí a necessidade da criação desta função.

\begin{lstlisting}[language=C, caption={Funções de libertação de memória.}]
/* -----------------------------------
 * Destructor Function Signatures
 * -----------------------------------
 */

void free_cfg_Lcampos( Lcampos campos );
void free_cfg_Conf( Conf cfg );
void free_cfg_Confs( Confs cfgs );
\end{lstlisting} 

Foi implementando a funções de libertação de memória para garantir que toda a memória, que inicialmente foi feita, o malloc, existe um free, garantindo assim que a aplicaçãos não tem qualquer tipo de memory leek.

\newpage
%%%%%%%%%%%%%%%%%%%%%%%%%%YACC-CONFIGUraÇÕES%%%%%%%%%%%%%%%%%%%%%%%%%%%%%
\subsubsection{YACC}
No ficheiro do \textit{YACC} existe, para cada produção, que seja "apanhada", uma função correspondente para tratar o caso apanhado.

\begin{lstlisting}[language=C, caption={YACC do ficheiro de configuração.}]
%token TITULO NPROVAS NTOP CAMPOS TEMPO CHAVE NOME str num
%{
#include <stdio.h>
#include <stdlib.h>
#include "cfg.lib.h"

#define YYERROR_VERBOSE

extern unsigned int cfglineno;
extern int cfglex (void);
extern int gco_set_Config(Confs cfg);
extern void cfglex_destroy();
%}
%union{
    char* cfg_str;
    int cfg_num;
    
    Confs cfg_confs;
    Conf cfg_conf;
    Lcampos cfg_campos;
}
%type <cfg_str> str Titulo
%type <cfg_num> num Nprovas Ntop Tempo Campo Chave Nome
%start Z

%type <cfg_confs> Z Configuracoes
%type <cfg_conf> Configuracao
%type <cfg_campos> Lcampos Campos
%%
Z : Configuracoes '$' { $$ = $1;
                        cfglex_destroy();
                        if( gco_set_Config($$) == 0 )
                            YYACCEPT;
                        YYABORT; };

Configuracoes : Configuracao '\n' Configuracoes {$$ = cons_cfg_Confs($1,$3);}
              | { $$ = cons_cfg_Confs_NIL(); }
              ;

Configuracao : Titulo  { $$ = cons_cfg_Conf_Titulo($1); }
             | Nprovas { $$ = cons_cfg_Conf_Nprovas($1); }
             | Ntop    { $$ = cons_cfg_Conf_Ntop($1); }
             | Tempo { $$ = cons_cfg_Conf_Tempo($1); }
             | Campos  { $$ = cons_cfg_Conf_Campos($1); }
             | Chave { $$ = cons_cfg_Conf_Chave($1); }
             | Nome { $$ = cons_cfg_Conf_Nome($1); }
             ;

Titulo  : TITULO str     { $$ = $2; };
Nprovas : NPROVAS num    { $$ = $2; };
Ntop    : NTOP num       { $$ = $2; };
Campos  : CAMPOS Lcampos { $$ = cfg_Lcampos_reverse( $2 ); }
Tempo   : TEMPO '$' num  { $$ = $3; };
Chave   : CHAVE '$' num  { $$ = $3; };
Nome    : NOME '$' num   { $$ = $3; };


Lcampos : Lcampos ';' Campo { $$ = cons_cfg_Lcampos_Lcampos($1,$3); }
        | Campo             { $$ = cons_cfg_Lcampos_Campo($1); }
        ;

Campo : '$' num { $$ = $2; };

%%
int yyerror( char* s ){
    fprintf(stderr, "Line %d: %s\n", cfglineno, s);
    return 0;
}
\end{lstlisting} 

Para o ficheiro de configuração o Yacc apanha os campos relativos à configuração do ficheiro. Cada uma das produções trata de um caso específico relativo aos campos. Assim, genericamente existe um campo e uma lista de campos, sendo necessário adicioná-lo à lista, daí exitir \texttt{\$\$ = cons\_cfg\_Lcampos\_Lcampos(\$1,\$3); }

A informação que o yacc recebe vem do processamento que o flex faz do ficheiro de configuração que ao flex é passado.
Quando a informação do flex é recebida pelo yacc, este apenas tem de a associar a uma das suas produções demodo a esta ser incluida na estrutura correta.


\newpage
%%%%%%%%%%%%%%%%%%%%%%%%%%LEX-CONFIGUraÇÕES%%%%%%%%%%%%%%%%%%%%%%%%%%%%%
\subsubsection{LEX}

\begin{lstlisting}[language=C, caption={Flex utilizado para o tratamento do ficherio de configuração}] 
%{
#define YY_NO_INPUT
#include <string.h>
#include "cfg.lib.h"
#include "cfg.tab.h"
%}
%option nounput
%option yylineno
%x ler_string
%x ler_numero
%x ler_tempo
%x campos
%%
"=titulo="  {BEGIN ler_string; return TITULO;}
"=nprovas=" {BEGIN ler_numero; return NPROVAS;}
"=ntop="    {BEGIN ler_numero; return NTOP;}
"=campos="  {BEGIN campos;     return CAMPOS;}
"=tempo="   {BEGIN ler_numero; return TEMPO;}
"=chave="   {BEGIN ler_numero; return CHAVE;}
"=nome="    {BEGIN ler_numero; return NOME;}

<ler_string>[^\n]+ {if(*cfgtext == ' ')
                        cfglval.cfg_str = strdup(cfgtext+1);
                    else
                        cfglval.cfg_str = strdup(cfgtext);
                    BEGIN 0;
                    return str;}
<ler_numero>\$     return '$';
<ler_numero>[0-9]+ {cfglval.cfg_num = atoi(cfgtext);
                    BEGIN 0;
                    return num;}

<campos>[0-9]+ { cfglval.cfg_num = atoi(cfgtext);
                 return num;}
<campos>\$     { return '$'; }
<campos>\;     { return ';'; }
<campos>\n     { BEGIN 0; return '\n'; }

<*><<EOF>>  return '$';

[\n]+       return '\n';

<*>[ \t] ;
<*>.     { /*printf("flex: %c (%d)\n", *cfgtext, *cfgtext);*/ }

%%

int cfgwrap(){
    return 1;
}
\end{lstlisting}

A informação que o ficheiro flex recebe vem do ficheiro de configuração.
Depois de reconhecermos os campos, este passam do flex para o yacc atrvés da estrtura que é interna no yacc.
Por exemplo, quando estamos a usar \texttt{cfglval.cfg\_str}, siginfica que estamos a passar (neste caso) uma string ao campo cfg\_str da union onde depois o yacc vai adiciona-lo à estrutura própria.

 
\newpage


%%%%%%%%%%%%%%%%%%%%%%%%%%CMD%%%%%%%%%%%%%%%%%%%%%%%%%%%%%
\subsection{Especificação para a Linguagem de Comandos da Consola}
A aplicação deverá ter um ambiente de consola com alguns comandos básicos:
\begin{itemize}
\item \textbf{Carregar configuração}: o sistema inicializa-se sempre com esta operação. Esta operação está sempre disponível para ser possível alterar a configuração.
\item \textbf{Carregar base de dados}: o sistema carrega para memória o estado de um torneio previamente gravado.
\item \textbf{Carregar resultado de prova}: o sistema carrega e processa um ficheiro CSV. Como resultado desta operação produz um ficheiro HTML com a pontuação de cada atleta, ordenado por ordem descente de pontos.
\item \textbf{Calcular ranking}: em qualquer momento poderá ser invocada esta operação à qual o sistema responder com uma listagem em HTML do raking;
\item \textbf{Gravar base de dados}: o sistema guarda em memória o estado da aplicação.
\item \textbf{Imprimir}: que apresenta a informação através de uma página HMTL.
\item \textbf{Save}: que guarda o estado atual do sistema.
\item \textbf{Sair}: o sistema termina a aplicação, se o estado atual não tiver sido previamente gravado deverá lançar um atleta e pedir configuração para prosseguir.
\end{itemize}


%%%%%%%%%%%%%%%%%%%%%%%%%%GRAMÁTICA-CMD%%%%%%%%%%%%%%%%%%%%%%%%%%%%%
\subsubsection{Especificação para a Linguagem de Comandos da Consola}
A gramática apresentar é a passada ao gabs para depois este gerar a estrutura usada para o funcionamento da aplicação em ambiente consola (CMD). Assim sendo tem-se:
\begin{lstlisting}[language=C, caption={Gramática para a linguagem de especificação para a linguagem de comandos de consola.}]
<ga>
  Cmds -> cons_Cmds(Cmd Cmds)
       |  cons_Cmds_NIL()
       ;

  Cmd -> cons_Cmd_Config(STR)
      |  cons_Cmd_Load(STR)
      |  cons_Cmd_Import(STR)
      |  cons_Cmd_Save(STR)
      |  cons_Cmd_FSave(STR)
      |  cons_Cmd_Print()
      |  cons_Cmd_Quit()
      |  cons_Cmd_FQuit()
      ;
</ga>
\end{lstlisting}

Esta gramática é definida depois de termos analisado quais as produções que podemos ter quando estamos a tratar os comandos.

\newpage

%%%%%%%%%%%%%%%%%%%%%%%%%%FUNÇÕES-CMD%%%%%%%%%%%%%%%%%%%%%%%%%%%%%
\subsubsection{Assinatura das funções utilizadas}
Vamos agora apresentar as assinaturas das funções que foram utilizadas para nós dar suporte à implementação do nosso projecto, nesta fase.

Quando falamos das funções temos: 
\begin{lstlisting}[language=C, caption={Construtores do ficheiro de consola.}]
.h do cmd
#ifndef __CMD_LIB_H
#define __CMD_LIB_H

#include <stdlib.h>

#define CMD_PROMPT "> "

// carrega uma nova configuracao
void cmd_config(char* ficheiro);

// carrega o estado da aplicacao
void cmd_load(char* ficheiro);

// importa um CSV
void cmd_import(char* ficheiro);

// escreve o html com os dados existentes
void cmd_print();

// guarda o estado da aplicacao
void cmd_save(char* ficheiro, int force);

// sair da aplicacao
int cmd_quit();

#endif
\end{lstlisting} 

Neste caso vale apena referir que quando temos as funções \texttt{void cmd\_config(char* ficheiro)}, e \texttt{void cmd\_import(char* ficheiro);}  estas quando em execução recebem o parser do ficheiro que lhe está associado, isto é, se temos o cmd\_config op parse que é invocado é o parser responsável pelo ficheiro de configuração e o cmd\_import "chama" o parse que é responsável pelo ficheiro CSV.

A função \texttt{void cmd\_print();} é a responsável por invocar a função que vai dar origem à página HTML pretendida.

\newpage


%%%%%%%%%%%%%%%%%%%%%%%%%%YACC-CMD%%%%%%%%%%%%%%%%%%%%%%%%%%%%%
\subsubsection{YACC}

No ficheiro do \textit{YACC} cada produção que seja "apanhada" é chamada a função correpondente para tratar a operação que foi capturada.

\begin{lstlisting}[language=C, caption={YACC do ficheiro cmd.}]
extern int cmdlex (void);
extern void cmdlex_destroy();
%}
%union{ char* cmd_ficheiro; }
%type <cmd_ficheiro> FICHEIRO
%start Comandos
%%
Comandos : Comando '\n' Comandos
         | 
         ;
Comando : Config
        | Load
        | Import
        | Print
        | Save
        | Quit
        | {/*| error  { printf("Comando invalido2.\n%s", CMD_PROMPT); yyclearin; yyerrok; }*/}
        ;

Config : CONFIG FICHEIRO { cmd_config($2); printf(CMD_PROMPT); };
       | CONFIG

Load : LOAD FICHEIRO { cmd_load($2); printf(CMD_PROMPT); };
     | LOAD

Import : IMPORT FICHEIRO { cmd_import($2); printf(CMD_PROMPT); }
       | IMPORT
       ;

Save : SAVE FICHEIRO { cmd_save($2, 0); printf(CMD_PROMPT); }
     | FSAVE FICHEIRO { cmd_save($2, 1); printf(CMD_PROMPT); }
     | SAVE
     | FSAVE
     ;

Print : PRINT { cmd_print(); printf(CMD_PROMPT); };

Quit : QUIT  { if(cmd_quit()){
                    cmdlex_destroy();
                    YYACCEPT;
               }else{
                    printf(CMD_PROMPT);
               }}
     | FQUIT { cmdlex_destroy(); YYACCEPT; }
     ;

%%
int yyerror( char* s ){
    //fprintf(stderr, "%s", s);
    return 0;
}
\end{lstlisting} 

Para o ficheiro de CMD tinhamos que ter o Yacc a apanhar os campos relativos aos comandos, para isso temos que cada produção trata um caso especifico relativo a esse campos.

\subsubsection{LEX}
\begin{lstlisting}[language=C, caption={LEX do ficheiro cmd.}]
%{
#define YY_NO_INPUT
#include <string.h>
#include "cmd.tab.h"
#include "cmd.lib.h"
%}
%option nounput
%x ficheiro
%%
"config " {BEGIN ficheiro; return CONFIG;}
"c "      {BEGIN ficheiro; return CONFIG;}

"load " {BEGIN ficheiro; return LOAD;}
"l "    {BEGIN ficheiro; return LOAD;}

"import " {BEGIN ficheiro; return IMPORT;}
"i "      {BEGIN ficheiro; return IMPORT;}

"print"  return PRINT;
"p"      return PRINT;

"save! " {BEGIN ficheiro; return FSAVE;}
"s! "    {BEGIN ficheiro; return FSAVE;}
"save "  {BEGIN ficheiro; return SAVE; }
"s "     {BEGIN ficheiro; return SAVE; }

"quit"   return QUIT;
"q"      return QUIT;
"quit!" return FQUIT;
"q!"    return FQUIT;

<<EOF>>  return FQUIT;

[\n]+    { return '\n'; }

<ficheiro>[^\n]+ { /* se houver erros, tentar um duplicate aqui */
                   //cmdlval.cmd_ficheiro = strdup(cmdtext);
                   cmdlval.cmd_ficheiro = cmdtext;
                   BEGIN 0;
                   return FICHEIRO;
                 }
<ficheiro>[\n]    { BEGIN 0;
                    printf("Comando invalido6.\n%s", CMD_PROMPT);
                    return '\n'; }

 /** capturar comandos invalidos especiais **/
[clis]\n          { printf("Comando invalido1.\n%s", CMD_PROMPT); }
[^clipsq\n][^\n]* { printf("Comando invalido2.\n%s", CMD_PROMPT); }
"q"[^\n]+         { printf("Comando invalido3.\n%s", CMD_PROMPT); }
"p"[^\n]+         { printf("Comando invalido4.\n%s", CMD_PROMPT); }
"config"[^ ]      { printf("Comando invalido5.\n%s", CMD_PROMPT); }

. ;

%%

void cmd_lex_push( FILE* sav_file_import ){
    cmdpush_buffer_state(cmd_create_buffer(sav_file_import, YY_BUF_SIZE));
}

void cmd_lex_pop(){
    yypop_buffer_state();
}

int cmdwrap(){
    return 1;
}
\end{lstlisting}

\subsection{Especificação para os Ficheiros de Resultados}
%%%%%%%%%%%%%%%%%%%%%%%%%%GRAMÁTICA-FICHEIROS-DE-RESULTADOS-%%%%%%%%%%%%%%%%%%%%%%%%%%%%%
\subsubsection{Gramática Especificação para os Ficheiros de Resultados}

A gramática a apresentar é a usada para o ficheiros de resultados. Assim sendo tem-se:
\begin{lstlisting}[language=C, caption={Gramática para a linguagem de especificação para os ficheiros de resultados.}]
<ga>
  
  Linhas -> cons_csv_Linhas(Linha Linhas)
          | cons_csv_Linhas_NIL()
          ;

  Linha -> cons_csv_Linha(Linha Campo)
        | cons_csv_Linha_Fim(Campo)
        ;

  Campo -> cons_csv_Campo(STR)
        | cons_csv_Campo_NIL()
        ;
 </ga>
\end{lstlisting}
\newpage
%%%%%%%%%%%%%%%%%%%%%%%%%%ESTRUTURA DE DADOS -FICHEIROS-DE-RESULTADOS-%%%%%%%%%%%%%%%%%%%%%%%%%%%%%
\subsubsection{Estrutura de Dados da Espeficicação dos Ficheiros de Resultados}

A estrutura que é gerada para o armazenamento dos campos que a gramática abstrata definiu, é divida consoante o nível em que estamos, isto é, nós começamos por ter a estrtura que vai abragir mais campos, neste caso o total de configurações, depois temos o que é um configuração e por último a estrutura que vai guardar um para uma configuração.

\begin{lstlisting}[language=C, caption={Estrutura de dados para armazenar a especificação dos ficheiros de resultados.}]

/* --- Linhas ---*/
struct sLinhas 
{ int flag;
  union {  
    struct {
        Linha s1;
        Linhas s2;
      } d1;
    struct {
      } d2;

  } u;
};

/* --- Linha ---*/
struct sLinha 
{ int flag;
  union {  
    struct {
        Linha s1;
        Campo s2;
      } d1;
    struct {
        Campo s1;
      } d2;

  } u;
};

/* --- Campo ---*/
struct sCampo 
{ int flag;
  union {  
    struct {
        char * s1;
      } d1;
    struct {
      } d2;

  } u;
};

\end{lstlisting}

\subsubsection{Assinatura das funções utilizadas}
Vamos agora apresentar as assinaturas das funções que foram utilizadas para nós dar suporte à implementação do nosso projecto, nesta fase.

Quando falamos das funções temos: 
\begin{lstlisting}[language=C, caption={Construtores do ficheiro de resultados.}]

/* -----------------------------------
 * Constructor Function Signatures
 * -----------------------------------
 */

Confs  cons_cfg_Confs( Conf a1, Confs a2);
Confs  cons_cfg_Confs_NIL();

Conf  cons_cfg_Conf_Titulo( char * a1);
Conf  cons_cfg_Conf_Nprovas( int a1);
Conf  cons_cfg_Conf_Ntop( int a1);
Conf  cons_cfg_Conf_Campos( Lcampos a1);
Conf  cons_cfg_Conf_Tempo( int a1);
Conf  cons_cfg_Conf_Chave( int a1);
Conf  cons_cfg_Conf_Nome( int a1);

Lcampos  cons_cfg_Lcampos_Lcampos( Lcampos a1, int a2);
Lcampos  cons_cfg_Lcampos_Campo( int a1);

/* -----------------------------------
 * Custom Function Signatures
 * -----------------------------------
 */

void cfg_Confs_print( Confs cfgs );
int cfg_Confs_validate( Confs cfgs );
Lcampos cfg_Lcampos_reverse( Lcampos l );

// getters
int cfg_get_Nprovas( Confs cfgs );
int cfg_get_Ntop( Confs cfgs );
int cfg_get_Tempo( Confs cfgs );
int cfg_get_Chave( Confs cfgs );
int cfg_get_Nome( Confs cfgs );
char* cfg_get_Titulo( Confs cfgs );

// obtem um array de inteiros com os campos. terminado com -1
int* cfg_get_Campos( Confs cfgs );
int cfg_get_NCampos( Confs cfgs );

// obtem um array de bytes que esta a 1 nos campos que foram seleccionados
char* cfg_Campos_seleccionado( Confs cfgs, int totalCampos );

/* -----------------------------------
 * Destructor Function Signatures
 * -----------------------------------
 */

void free_cfg_Lcampos( Lcampos campos );
void free_cfg_Conf( Conf cfg );
void free_cfg_Confs( Confs cfgs );

#endif
\end{lstlisting} 
\newpage

\subsubsection{YACC}
%%%%%%%%%%%%%%%%%%%%%%%%%%YACC-FICHEIROS-DE-RESULTADOS-%%%%%%%%%%%%%%%%%%%%%%%%%%%%%
No ficheiro do \textit{YACC} cada produção que seja "apanhada" é chamada a função correpondente para tratar o caso apanhado.

\begin{lstlisting}[language=C, caption={YACC do ficheiro de resultados.}]

%token str
%{
#include <stdio.h>
#include <stdlib.h>
#include "csv.lib.h"

#define YYERROR_VERBOSE

extern unsigned int csvlineno;
extern int csvlex (void);
extern void csvlex_destroy();
%}
%union{
  char * csv_str;
  
  Linhas csv_Linhas;
  Linha  csv_Linha;
  Campo  csv_Campo;

} 

%type<csv_str> str
%type<csv_Campo> Campo
%type<csv_Linhas> X Linhascsv
%type<csv_Linha> Linha
%start X

%%
X : Linhascsv '$' { $$ = $1; 
                    csvlex_destroy();
                    if( csv_Linhas_validate($$) ){
                        csv_import_csv($$);
                        YYACCEPT;
                    }else{
                        YYABORT;
                    }
                  };

Linhascsv : Linha '\n' Linhascsv {$1 = csv_Linha_reverse($1);
                                  $$ = cons_csv_Linhas ($1,$3);}
          | {$$ = cons_csv_Linhas_NIL();}
          ;

Linha : Linha ';' Campo {$$ = cons_csv_Linha ($1,$3);}
      | Campo {$$ = cons_csv_Linha_Fim ($1);}
      ;

Campo : str {$$ = cons_csv_Campo ($1);}
      | {$$ = cons_csv_Campo_NIL();}
      ;

%%
int yyerror( char* s ){
    fprintf(stderr, "Line %d: %s\n", csvlineno, s);
    return 0;
}
\end{lstlisting}
%%%%%%%%%%%%%%%%%%%%%%%%%%LEX-FICHEIROS-DE-RESULTADOS-%%%%%%%%%%%%%%%%%%%%%%%%%%%%%
\subsubsection{LEX}

\begin{lstlisting}
%{
#define YY_NO_INPUT
#include <string.h>
#include "csv.lib.h"
#include "csv.tab.h"
%}
%option yylineno
%option nounput

%%
[^\n\;]+ {csvlval.csv_str = strdup(csvtext);
         return str;}
\n+     {return '\n';}
\;      {return ';';}
<<EOF>> {return '$';}


%%
int csvwrap(){
  return 1;
}
\end{lstlisting}
\newpage
\subsection{Tratamento da Informação do Atleta}

Um dos requesitos necessários para o desenvolvimento da plataforma era o cálculo da pontuação dos atletas. 
Assim, foram implementadas funções que permitem determinar os rankings necessários.
\begin{lstlisting}[language=C, caption={Estrutura de um atleta.}]
struct sLInt {
    LInt proximo;
    int num;
};

struct sAtletas{
    Atletas proximo;
    Atleta atleta;
};

struct sAtleta{
    char* ID;
    char* nome;
    LInt scores;
    int tempo_ultima_prova;
};
\end{lstlisting}

\begin{lstlisting}[language=C, caption={Assinatura das funções.}]
Atleta atl_novo_atleta(char* id, char* nome);

void atl_add_score(Atleta atleta, int score);
int atl_get_score(Atleta atleta, int nTop);

Atletas atl_ordenar_por_Score(Atletas atletas);
Atletas atl_insere_por_ID(Atletas atletas, Atleta atleta);
Atletas atl_insere_por_Score(Atletas atletas, Atleta atleta);

void atl_print(Atletas atletas);

void atl_free_Atletas(Atletas atletas);

// versao que nao liberta memoria dos atletas em si (apenas da lista)
// usado para libertar a lista ordenada por score
void atl_free_AtletasList(Atletas atletas);
void atl_free_Atleta(Atleta atleta);
void atl_free_Scores(LInt scores);


/******************
 * Ler dados dos atletas a partir do csv
 * ***************/

Atletas atl_ler_csv( Confs cfg, Linhas csv, Atletas atletas );
Atletas atl_ler_tempos( Confs cfg, Linhas csv, Atletas atletas );

#endif
\end{lstlisting}

\newpage
%%%%%%%%%%%%%%%%%%%%%%%%%HTML%%%%%%%%%%%%%%%%%%%%%%%%%%
\subsection{HTML}

Um dos resultados finais deste projeto é o de gerar uma página HTML com a pontuação de cada atleta ordenado por ordem decrescente de pontos de pontos. 

\begin{lstlisting}[language=C, caption={Assinatura das funções.}]
void htm_print_indice();
void htm_print_provas();
void htm_print_resultados();
int htm_begin();
void htm_end();
\end{lstlisting}

\newpage
\section{Conclusão}

\newpage
\section{Anexos}

\subsection{Makefile}

\begin{lstlisting}[language=C, caption={Makefile }]
.PHONY: clean

GDB=-Wall -g

#----------------\
#  Tokenizers    \
#----------------\

cmdtok : cmd.tok.c lex.cmd.o cmd.tab.o cmd.lib.o cmd.mockups.o
	gcc $(GDB) -o $@ $^ -ll

cfgtok : cfg.tok.c lex.cfg.o cfg.tab.o cfg.lib.o
	gcc $(GDB) -o $@ $^ -ll

csvtok : csv.tok.c lex.csv.o csv.tab.o csv.lib.o
	gcc $(GDB) -o $@ $^ -ll

#----------------\
#  Gco           \
#----------------\

gco : gco.lib.c sav.lib.o htm.lib.o atl.lib.o cfg.lib.o lex.cfg.o cfg.tab.o cmd.lib.o lex.cmd.o cmd.tab.o csv.lib.o lex.csv.o csv.tab.o
	gcc $(GDB) -o $@ $^
	cp gco ../bin/gco

htm.lib.o : htm.lib.c htm.lib.h cfg.lib.h csv.lib.h
	gcc $(GDB) -c htm.lib.c

atl.lib.o : atl.lib.c atl.lib.h cfg.lib.h csv.lib.h
	gcc $(GDB) -c atl.lib.c

sav.lib.o : sav.lib.c sav.lib.h
	gcc $(GDB) -c sav.lib.c

#----------------\
#  Cfg           \
#----------------\

cfg : cfg.main.c cfg.lib.o lex.cfg.o cfg.tab.o
	gcc $(GDB) -o $@ $^ -ll

cfg.lib.o : cfg.lib.c cfg.lib.h
	gcc $(GDB) -c cfg.lib.c

cfg.tab.o : cfg.tab.h cfg.tab.c
	gcc $(GDB) -c cfg.tab.c

lex.cfg.o : lex.cfg.c cfg.tab.h cfg.lib.o
	gcc $(GDB) -c lex.cfg.c -ll

cfg.tab.c cfg.tab.h : cfg.y cfg.lib.o
	yacc -d -pcfg -bcfg cfg.y

lex.cfg.c : cfg.l
	flex -Pcfg cfg.l

#----------------\
#  Cmd           \
#----------------\

cmd : cmd.main.c cmd.lib.o lex.cmd.o cmd.tab.o cmd.mockups.o
	gcc $(GDB) -o $@ $^ -ll

cmd.mockups.o : cmd.mockups.c cmd.mockups.h
	gcc $(GDB) -c cmd.mockups.c

cmd.lib.o : cmd.lib.c cmd.lib.h
	gcc $(GDB) -c cmd.lib.c

cmd.tab.o : cmd.tab.h cmd.tab.c
	gcc $(GDB) -c cmd.tab.c

lex.cmd.o : lex.cmd.c cmd.tab.h
	gcc $(GDB) -c lex.cmd.c -ll

cmd.tab.c cmd.tab.h : cmd.y cmd.lib.h
	yacc -d -pcmd -bcmd cmd.y

lex.cmd.c : cmd.l
	flex -Pcmd cmd.l


#----------------\
#  CSV           \
#----------------\

csv : csv.main.c csv.lib.o lex.csv.o csv.tab.o
	gcc $(GDB) -o $@ $^ -ll

csv.lib.o : csv.lib.c csv.lib.h
	gcc $(GDB) -c csv.lib.c

csv.tab.o : csv.tab.h csv.tab.c
	gcc $(GDB) -c csv.tab.c

lex.csv.o : lex.csv.c csv.tab.h
	gcc $(GDB) -c lex.csv.c -ll

csv.tab.c csv.tab.h : csv.y csv.lib.h
	yacc -d -pcsv -bcsv csv.y

lex.csv.c : csv.l
	flex -Pcsv csv.l
#----------------\
#  Extras        \
#----------------\

clean :
	$(RM) gabs
	$(RM) gco
	$(RM) cmd
	$(RM) cfg
	$(RM) csv
	$(RM) *.o
	$(RM) *.tab.[ch]
	$(RM) lex.*.c
	$(RM) *tok

gabs:
	cd ./../gabs/ && $(MAKE) gabs && mv ./gabs ./../src/gabs

\end{lstlisting}
\newpage

\section{Elementos do Grupo}
\begin{figure}[h!]
\centering
\begin{subfigure}{.33\textwidth}
  \centering
  \includegraphics[width=0.8\linewidth]{60}
  \caption{Bruno Ferreira  }
\end{subfigure}%
\begin{subfigure}{.33\textwidth}
  \centering
  \includegraphics[width=0.8\linewidth]{107}
  \caption{Cláudia Oliveira}
\end{subfigure}%
\begin{subfigure}{.33\textwidth}
  \centering
  \includegraphics[width=0.8\linewidth]{93}
  \caption{Vanessa Campos}
\end{subfigure}%
\end{figure}



\end{document}